Algorithmic redistricting has been developed almost exclusively in the American single-member-district context. We apply edge-weighted recursive bisection to six parliamentary democracies with distinct electoral systems: the United Kingdom (FPTP, 543 single-member constituencies), Canada (FPTP, 343 ridings), Ireland (STV, 43 multi-member constituencies), Malta (STV, 13 constituencies with 5 seats each), New Zealand (MMP, 65 general electorates), and Germany (MMP, 299 Wahlkreise). For each country, we use the RPLAN open interchange format to draw algorithmically optimal boundaries and compare compactness and population balance against enacted boundaries. We find that algorithmic redistricting transfers successfully across electoral systems, achieving 10--22\% compactness improvements over enacted boundaries in all six countries. For multi-member systems (Ireland, Malta), we introduce a \emph{per-seat balance} criterion and demonstrate that the D'Hondt seat allocation method produces near-perfect proportionality (Gallagher index $< 2.0$) even when boundaries are drawn purely on geographic grounds. We release the first open-source international redistricting dataset covering 1,333 constituencies across six countries, enabling cross-national comparison of redistricting algorithms and outcomes.
